\documentclass[a4paper]{report}
\usepackage[fontsize=14pt]{scrextend}
\usepackage[
top=0.85in,
bottom=1in,
inner=0.675in,
outer=0.65in,
]{geometry}
\usepackage{setspace}
\setstretch{1.5}

\usepackage{needspace}
\usepackage{etoolbox}

\usepackage{enumitem}
\setlist{nosep}

\usepackage{graphicx}
\usepackage{xcolor}
\usepackage{etoc}


\usepackage{amsmath, amssymb, amsthm}
\usepackage{autoparen}
\usepackage{thmtools}

\definecolor{thmred}{HTML}{e86d70}
\definecolor{defgreen}{HTML}{14b840}
\definecolor{exorange}{HTML}{a88c00}

\usepackage{tikz}
\usetikzlibrary{decorations.markings, arrows.meta, calc, angles, quotes}

\tikzset{
	tri/.tip={Triangle[angle=30:8pt]},
	dashed/.style={dash pattern=on 5pt off 3pt},
	arrow at/.style n args={3}{
		postaction={decorate,
			decoration={markings,
				mark=at position #1 with {
					\arrow[xshift=5.25pt]{tri};
					\node[#3] {#2};
				}
			}
		}
	},
	arrow at/.default=0.5,
}
% ---- measureline: FIXED VERSION (no \p1/\n1 — macro-safe) ----
\newcommand{\measureShrink}{0.34}

% \measureline{start}{end}{label}{label pos}{offset}
\newcommand{\measureline}[5]{%
	\begin{scope}
		\coordinate (mlS) at (#1);
		\coordinate (mlE) at (#2);
		
		\pgfpointdiff{\pgfpointanchor{mlS}{center}}{\pgfpointanchor{mlE}{center}}
		\pgfgetlastxy{\mlDx}{\mlDy}
		\pgfmathsetmacro{\mlAng}{atan2(\mlDy,\mlDx)}
		
		\coordinate (mlSin) at ($(mlS)+(\mlAng+90:#5)+(\mlAng:\measureShrink)$);
		\coordinate (mlEin) at ($(mlE)+(\mlAng+90:#5)-(\mlAng:\measureShrink)$);
		
		
		
		\draw[thick, {Bar[width=6pt]}-{Bar[width=6pt]}] (mlSin) -- (mlEin)
		node[midway, #4] {#3};
	\end{scope}
}

\newcounter{myfig}
\newenvironment{myfigure}[1][]{%
	\par\vspace{\intextsep}\noindent
	\begin{minipage}{\linewidth}
		\def\mycaption{#1}%
	}{%
		\refstepcounter{myfig}%
		\quad\noindent
		\ifx\mycaption\empty
		\textbf{الشكل \themyfig}%
		\else
		\begin{tabular}[b]{@{}l@{}}
			\textbf{الشكل \themyfig} \\[-5pt]
			\mycaption
		\end{tabular}%
		\fi
	\end{minipage}%
	\par\vspace{\intextsep}%
}

\usepackage{environ}
\NewEnviron{exercises}{%
	\section*{تمارين}%
	\begin{tasks}[cols=1, label=\textbf{\arabic*.}]
		\BODY
	\end{tasks}
} 


\AtBeginDocument{
	\allowdisplaybreaks
	\renewcommand{\jot}{10pt}
	\abovedisplayskip=5pt
	\belowdisplayskip=10pt
}



\usepackage[main=arabic, bidi=basic, provide=*, layout=graphics]{babel}
\babelprovide[import,justification=kashida, transforms=kashida.base]{arabic}
\babelfont{rm}[Script=Arabic, ItalicFont=*, ItalicFeatures={FakeSlant=0.2}]{Times New Roman}

\babelfont[english]{rm}{Times New Roman}
\newcommand{\en}[1]{\foreignlanguage{english}{#1}}

\newfontfamily{\farsi}{Gulzar}[Script=Arabic]
\newfontfamily{\kufi}{Reem Kufi}[Script=Arabic]



\usepackage[lite, zswash, subscriptcorrection, straightbraces]{mtpro2}
\DeclareMathSymbol{0}{\mathalpha}{operators}{`0}
\DeclareMathSymbol{1}{\mathalpha}{operators}{`1}
\DeclareMathSymbol{2}{\mathalpha}{operators}{`2}
\DeclareMathSymbol{3}{\mathalpha}{operators}{`3}
\DeclareMathSymbol{4}{\mathalpha}{operators}{`4}
\DeclareMathSymbol{5}{\mathalpha}{operators}{`5}
\DeclareMathSymbol{6}{\mathalpha}{operators}{`6}
\DeclareMathSymbol{7}{\mathalpha}{operators}{`7}
\DeclareMathSymbol{8}{\mathalpha}{operators}{`8}
\DeclareMathSymbol{9}{\mathalpha}{operators}{`9}

\newcommand{\arabicordinal}[1]{%
	\ifcase#1
	\or الأول
	\or الثاني
	\or الثالث
	\or الرابع
	\or الخامس 
	\or السادس
	\or السابع
	\or الثامن
	\or التاسع
	\or العاشر
	\else \number#1
	\fi
}

\renewcommand{\thechapter}{\arabicordinal{\value{chapter}}}

\etocsettocstyle{
	\vspace*{10pt}\noindent
	{\farsi\fontsize{25pt}{20pt}\textbf{المحتويات}}
	\vspace*{10pt}
}{}

\etocsetlevel{chapter}{-1}

\etocsetstyle{chapter}{}{\par}{
	\vspace{15pt}\noindent\textbf{الفصل \etocnumber : \etocname}
}

\etocsetlevel{schapter}{0}
\etocsetstyle{schapter}{}{\par}{\vspace{15pt}\noindent\textbf{\etocname}\hfill\etocpage}


\etocsetlevel{section}{1}
\etocsetstyle{section}{}{\par}{\noindent\etocnumber~~\etocname\hfill\etocpage}


\etocsetlevel{subsection}{2}
\etocsetstyle{subsection}{}{\par}{\noindent\etocnumber~~\etocname\hfill\etocpage}

\usepackage{enumitem}
\usepackage[spacing=1.3]{tasksenv}
\makeatletter





\def\@makechapterhead#1{%
	\vspace*{10pt}
	{\noindent\farsi\fontsize{35pt}{35pt} \textbf{الفصل \thechapter} \\[25pt]
	\kufi\textbf{#1}
	}
	\vspace*{150pt}
	\thispagestyle{empty}
}

\def\@makeschapterhead#1{%
	\vspace*{5\p@}%
	{\parindent 0pt \centering
		\normalfont
		\interlinepenalty1000
		\large \bfseries  #1\par\nobreak
		\vskip 25\p@
}}


\renewenvironment{proof}[1]{
	\par\noindent
	\colorbox{gray!15}{\textbf{البرهان}}\par
	\noindent#1}{$\blacksquare$\par}

\makeatother

\newenvironment{solution}[1]{
	\par\noindent
	\colorbox{gray!15}{\textbf{الحل}}\par
	\noindent#1}{}

\newenvironment{note}[1]{
	\par\noindent
	\colorbox{gray!15}{\textbf{ملاحظة}}\par
	\noindent#1}{\par}

\renewcommand{\thesection}{\arabic{section}}

% ---- per-section total tracking (2-pass, via the .aux file) -------------
\makeatletter
\newcommand{\@setsectiontotal}[3]{%
	\expandafter\gdef\csname sectiontotal@#1@#2\endcsname{#3}}
\newcommand{\getsectiontotal}[2]{%
	\ifcsname sectiontotal@#1@#2\endcsname
		\csname sectiontotal@#1@#2\endcsname
	\else
		2 % unknown on the first pass: default to "more than one"
	\fi
}
\newcommand{\sectionkey}{\arabic{chapter}-\arabic{section}}
\newcommand{\thmheadnumber}[1]{%
	\immediate\write\@auxout{%
		\string\@setsectiontotal{#1}{\sectionkey}{\the\value{#1}}}%
	\ifnum\getsectiontotal{#1}{\sectionkey}>1 ~\NUMBER\fi
}
\makeatother

% Theorem styles
\declaretheoremstyle[
	headfont=\bfseries,
	headformat={\colorbox{thmred}{\NAME\thmheadnumber{theorem}}\\},
	headpunct={},
	bodyfont=\itshape
]{thmstyle}

\declaretheoremstyle[
	headfont=\bfseries,
	headformat={\colorbox{exorange}{\NAME\thmheadnumber{example}}\\},
	headpunct={},
	bodyfont=\normalfont
]{exstyle}

\declaretheoremstyle[
	headfont=\bfseries,
	headformat={\colorbox{defgreen}{\NAME\thmheadnumber{definition}}\\},
	headpunct={},
	bodyfont=\normalfont
]{defstyle}

% Theorem environments
\declaretheorem[style=thmstyle, name=مبرهنة, numberwithin=section]{theorem}
\declaretheorem[style=exstyle, name=مثال, numberwithin=section]{example}
\declaretheorem[style=defstyle, name=تعريف, numberwithin=section]{definition}
\renewcommand{\thetheorem}{\arabic{theorem}}
\renewcommand{\theexample}{\arabic{example}}
\renewcommand{\thedefinition}{\arabic{definition}}


\newcommand{\N}{\mathbb{N}}
\newcommand{\Z}{\mathbb{Z}}
\newcommand{\Q}{\mathbb{Q}}
\newcommand{\R}{\mathbb{R}}
\newcommand{\C}{\mathbb{C}}

\newcommand{\HW}{\textcolor{red}{\en{H.W.}}}

\DeclareMathOperator{\re}{Re}
\DeclareMathOperator{\im}{Im}
\DeclareMathOperator{\Arg}{Arg}
\preto{\section}{\Needspace{4\baselineskip}}
\preto{\subsection}{\Needspace{4\baselineskip}}

\AtBeginEnvironment{theorem}{\Needspace{4\baselineskip}}
\AtBeginEnvironment{definition}{\Needspace{4\baselineskip}}
\AtBeginEnvironment{example}{\Needspace{4\baselineskip}}
\AtBeginEnvironment{proof}{\Needspace{4\baselineskip}}
\AtBeginEnvironment{solution}{\Needspace{4\baselineskip}}

\usepackage{mathshorthand}
